\begin{figure}[tb]
\centering
\begin{tikzpicture}[scale=1.05,>=Latex,font=\small]
\coordinate (O) at (0,0);
\draw[->,gray] (O)--(3.1,0) node[right]{\(d\)};
\draw[->,gray] (O)--(0,2.5) node[above]{\(q\)};
\draw[->,MidnightBlue,very thick] (O)--({2.7*cos(32)},{2.7*sin(32)}) node[right]{\(\gamma\)-axis};
\draw[->,MidnightBlue,thick] (O)--({2.0*cos(122)},{2.0*sin(122)}) node[above]{quadrature};
\draw[->,Purple,very thick] (O)--({2.25*cos(60)},{2.25*sin(60)}) node[above]{\(\vect x\)};
\draw[->] (0.75,0) arc[start angle=0,end angle=32,radius=.75];
\node at (0.82,.23) {\(\delta\)};
\node[draw,rounded corners,align=left,fill=Orange!10] at (6.3,1.2)
 {\(\delta=\mathrm{const.}\): rotor-fixed\\
  \(\delta=\delta(\vect i,\boldsymbol\psi)\): moving\\
  extra rate: \(-\dot\delta\matr J\vect x_\gamma\)};
\end{tikzpicture}
\caption{An arbitrary rotor-relative spatial frame. A state-dependent angle
simplifies instantaneous geometry but contributes its own angular velocity.}
\label{fig:arbitrary-frame}
\end{figure}
